\documentclass[a4paper,12pt]{article}
\usepackage[utf8]{inputenc}
\usepackage[russian]{babel}
\usepackage{graphicx}
\usepackage{hyperref}
\usepackage{minted}
\usemintedstyle{friendly}
\usepackage{xcolor}
\usepackage{geometry}
\usepackage{amsmath}
\usepackage{enumitem}
\usepackage{float}
\usepackage{indentfirst}
\usepackage{caption}
\usepackage{amsfonts}
\usepackage{tikz}
\usetikzlibrary{shapes.geometric, arrows.meta, positioning}

\documentclass[14pt]{article}
\usepackage[utf8]{inputenc}
\usepackage[russian]{babel}
\usepackage{geometry}
\usepackage{minted}

\geometry{top=2cm, bottom=2cm, left=2.5cm, right=1.5cm}

% Настройки для листингов кода
\setminted{
    fontsize=\small,
    breaklines=true,
    autogobble=true,
    frame=single,
}

\begin{document}

\begin{titlepage}
    \centering
    % Верхний блок
    {\bfseries МИНИСТЕРСТВО НАУКИ И ВЫСШЕГО ОБРАЗОВАНИЯ РФ \\
    ФГБОУ ВО «Белгородский государственный технологический университет им. В.Г. Шухова»\par}
    
    \vfill
    
    % Центральный блок
    {\bfseries\Large Практическое задание №8\par}
    \vspace{0.5cm}
    {\large по дисциплине: Основы программирования и алгоритмизации \par}
    {\large тема: «Геометрия и теория множеств» \par}
    
    \vfill
    
    % Блок авторов
    \begin{flushright}
    \begin{tabular}{l}
        \large Выполнил: студент группы ВТ-252 \\
        \large Дронь Дмитрий Алексеевич \\
        \large Проверила: \\
        \large Стеценко Оксана Николаевна
    \end{tabular}
    \end{flushright}
    
    \vfill
    
    % Нижний блок
    {\large Белгород, 2026 г. \par}
\end{titlepage}

\newpage
\begin{center}
\section*{Цель работы}
\end{center}
Цель методических указаний организовать самостоятельную работу студента по выполнению индивидуальных заданий, направленных на закрепление навыков программирования на языке высокого уровня и способов решения математических задач с применением программы для работы с электронными таблицами MS Excel.
\newpage
\section*{Индивидуальное задание №8}
\textbf{Задача 5.} Найти расстояние от заданной точки до заданного треугольника на плоскости.
\newpage
\begin{center}
\section*{Блок-схема}
\end{center}
\begin{center}
\scalebox{0.90}{
\begin{tikzpicture}[
    >=Stealth,
    line/.style={draw,-Stealth,thick},
    block/.style={
        rectangle,
        rounded corners,
        draw,
        align=center,
        text width=3.8cm,
        minimum height=1cm
    },
    blockplain/.style={
        rectangle,
        draw,
        align=center,
        text width=3.8cm,
        minimum height=1cm
    },
    decision/.style={
        diamond,
        draw,
        align=center,
        aspect=2.2,
        text width=2.6cm,
        inner sep=2pt
    }
]

%------------------ Узлы ------------------
\node[block] (input) at (0,0)
{
dist\_point\_segment\\
px, py, ax, ay, bx, by
};

\node[blockplain] (init) at (0,-3.5)
{
x = bx - ax\\
vy = by - ay\\
wx = px - ax\\
wy = py - ay\\
c1 = dot(wx, wy, vw, vy)\\
c2, dist
};

\node[decision] (c1check) at (0,-7.5)
{c1 <= 0};

% Ветка да (вправо)
\node[blockplain] (distA) at (4.5,-10)
{
dist = sqrt(norm2(px - ax, py - ay))
};

% Ветка нет (влево)
\node[blockplain] (c2calc) at (-4.5,-10)
{
c2 = dot(vx, vy, vx, vy)
};

\node[decision] (c2check) at (-4.5,-13.5)
{c2 <= c1};

% c2 <= c1: да (вправо)
\node[blockplain] (distB) at (0,-16.5)
{
dist = sqrt(norm2(px - bx, py - by))
};

% c2 <= c1: нет (влево)
\node[blockplain] (proj) at (-9,-16.5)
{
t = c1 / c2;\\
projx = ax + t * vx\\
projy = ay + t * vy\\
dist = sqrt(norm2(px - projx, py - projy))
};

\node[block] (ret) at (0,-22.5)
{возврат dist};

%------------------ Стрелки ------------------
\draw (input) -- (init);
\draw (init) -- (c1check);

% c1check: да -> distA (вправо)
\draw
    (c1check.east) -- ++(2.65,0)
    node[midway,above]{да}
    -- (distA.north);

% c1check: нет -> c2calc (влево)
\draw
    (c1check.west) -- ++(-2.65,0)
    node[midway,above]{нет}
    -- (c2calc.north);

% c2calc -> c2check
\draw (c2calc.south) -- (c2check.north);

% c2check: да -> distB (вправо)
\draw
    (c2check.east) -- ++(2.65,0)
    node[midway,above]{да}
    -- (distB.north);

% c2check: нет -> proj (вниз)
\draw
    (c2check.west) -- ++(-2.65,0)
    node[midway,above]{нет}
    -- (proj.north);

% distA -> ret (вниз и влево)
\draw
    (distA.south) -- ++(0,-10.28)
    -- ++(-4.5,0)
    -- (ret.north);

% distB -> ret (вниз)
\draw
    (distB.south) -- ++(0,-2.2775)
    -- ++(-4.5,0)
    -- ++(0,-1.5)
    -- ++(4.5,0)
    -- (ret.north);

% proj -> ret (вниз и вправо)
\draw
    (proj.south) -- ++(0,-1.5)
    -- ++(4.5,0)
    -- ++(0,-1.5)
    -- ++(4.5,0)
    -- (ret.north);

\end{tikzpicture}
}
\end{center}
\newpage
\begin{center}
\scalebox{0.90}{
\begin{tikzpicture}[
    >=Stealth,
    line/.style={draw,-Stealth,thick},
    block/.style={
        rectangle,
        rounded corners,
        draw,
        align=center,
        text width=4.5cm,
        minimum height=1cm
    },
    blockplain/.style={
        rectangle,
        draw,
        align=center,
        text width=4.5cm,
        minimum height=1cm
    },
    decision/.style={
        diamond,
        draw,
        align=center,
        aspect=2.2,
        text width=2.8cm,
        inner sep=2pt
    }
]

%------------------ Узлы ------------------
\node[block] (input) at (0,0)
{
compute\_distance\\
px, py, ax, ay, bx, by,\\
cx, cy\\
r.error\_flag = 0\\
r.dist = 0.0
};

\node[blockplain] (calc) at (0,-5)
{
d1 = dist\_point\_segment(px, py, ax, ay, bx, by)\\
d2 = dist\_point\_segment(px, py, bx, by, cx, cy)\\
d3 = dist\_point\_segment(px, py, cx, cy, ax, ay)\\
min = d1
};

\node[decision] (d2check) at (0,-9.5)
{d2 < min};

\node[blockplain] (mind2) at (0,-12)
{min = d2};

\node[decision] (d3check) at (0,-14.5)
{d3 < min};

\node[blockplain] (mind3) at (0,-17)
{min = d3};

\node[blockplain] (rdist) at (0,-19.5)
{r.dist = min};

\node[block] (ret) at (0,-22)
{возврат r};

%------------------ Стрелки ------------------
\draw (input) -- (calc);
\draw (calc) -- (d2check);

% d2check: mind2
\draw
    (d2check.south) -- ++(0,-0.2)
    -- (mind2.north);

\draw (mind2.south) -- (d3check.north);

% d3check: mind3
\draw
    (d3check.south) -- ++(0,-0.2)
    -- (mind3.north);



\draw (mind3.south) -- (rdist.north);
\draw (rdist.south) -- (ret.north);

\end{tikzpicture}
}
\end{center}

\newpage
\section*{Программа}
\begin{minted}{c}
/* Летняя практика. Индивидуальное задание №8 «Геометрия и теория множеств»
   Выполнил студент группы ВТ-252 Дронь Д.А.
   Задание 5: Найти расстояние от заданной точки до заданного треугольника на плоскости.
*/
#include <stdio.h>
#include <math.h>

// Структура результата
typedef struct {
    double dist;
    int error_flag;
} Result;

// Ввод координат
int input_data(double *px, double *py,
               double *ax, double *ay,
               double *bx, double *by,
               double *cx, double *cy)
{
    int flag = 0;

    printf("Введите координаты точки P(x y): ");
    if (scanf("%lf %lf", px, py) != 2) {
        flag = 1;
    }

    if (flag == 0) {
        printf("Введите координаты A(x y), B(x y), C(x y): ");
        if (scanf("%lf %lf %lf %lf %lf %lf",
                  ax, ay, bx, by, cx, cy) != 6) {
            flag = 1;
        }
    }

    return flag;
}

// Скалярное произведение
double dot(double x1, double y1, double x2, double y2) {
    return x1 * x2 + y1 * y2;
}

// Квадрат длины вектора
double norm2(double x, double y) {
    return x * x + y * y;
}

/*
Расстояние от точки P до отрезка AB
(если проекция вне отрезка — возвращаем расстояние до ближайшей вершины)
*/
double dist_point_segment(double px, double py,
                          double ax, double ay,
                          double bx, double by)
{
    double vx = bx - ax;
    double vy = by - ay;
    double wx = px - ax;
    double wy = py - ay;

    double c1 = dot(wx, wy, vx, vy);
    if (c1 <= 0) {
        return sqrt(norm2(px - ax, py - ay));
    }

    double c2 = dot(vx, vy, vx, vy);
    if (c2 <= c1) {
        return sqrt(norm2(px - bx, py - by));
    }

    double t = c1 / c2;
    double projx = ax + t * vx;
    double projy = ay + t * vy;

    return sqrt(norm2(px - projx, py - projy));
}

// Основная функция вычисления расстояния
Result compute_distance(double px, double py,
                        double ax, double ay,
                        double bx, double by,
                        double cx, double cy)
{
    Result r;
    r.error_flag = 0;
    r.dist = 0.0;

    double d1 = dist_point_segment(px, py, ax, ay, bx, by);
    double d2 = dist_point_segment(px, py, bx, by, cx, cy);
    double d3 = dist_point_segment(px, py, cx, cy, ax, ay);

    double min = d1;
    if (d2 < min) {
        min = d2;
    }
    if (d3 < min) {
        min = d3;
    }

    r.dist = min;
    return r;
}

// Вывод результата
void output_result(Result r) {
    printf("Расстояние от точки до треугольника: %.6f\n", r.dist);
}

int main(void) {
    double px, py, ax, ay, bx, by, cx, cy;

    int err = input_data(&px, &py, &ax, &ay, &bx, &by, &cx, &cy);

    switch (err) {
        case 0:
            break;
        case 1:
            printf("Ошибка: некорректный формат ввода.\n");
            return 1;
        default:
            printf("Неизвестная ошибка ввода.\n");
            return 1;
    }

    Result r = compute_distance(px, py, ax, ay, bx, by, cx, cy);

    switch (r.error_flag) {
        case 0:
            output_result(r);
            break;
        default:
            printf("Неизвестная ошибка вычислений.\n");
            break;
    }

    return 0;
}


\end{minted}
\newpage
\section*{Тесты}
\begin{figure}[H]
    \centering
    \includegraphics[width=0.9\textwidth]{test.png}
    \caption{Результат выполнения программы}
\end{figure}
\newpage
\section*{Вывод}
В ходе выполнения работы мне удалось организовать самостоятельную работу по выполнению индивидуальных заданий, направленных на закрепление навыков программирования на языке высокого уровня и способов решения математических задач с применением программы для работы с электронными таблицами MS Excel.

\end{document}




